\documentclass[../../../../main.tex]{subfiles}
\begin{document}

\begin{flushright}
\footnotesize\textit{【自動文字起こし・要確認】}
\end{flushright}

[解] $-\frac{\pi}{2} \leqq x \leqq \frac{\pi}{2}$ で考える。以下 $C = \cos t, S = \sin t$ とする。

まずAについて、面積が最大になるのは、明らかに4頂点が $(\pm t, 0), (\pm t, \cos t)$ で与えられる時で ($0 < t < \pi/2$) この時
\[
A = 2t C \quad (0 < t < \pi/2)
\]
である。
\[
\frac{dA}{dt} = 2(C - t S) = 2C (1 - t \tan t)
\]
より、下表をとる。($t_0$ は $t_0 \tan t_0 = 1$ を満たす)

\begin{center}
\begin{tabular}{c|c|c|c|c|c}
$t$ & 0 & $\cdots$ & $t_0$ & $\cdots$ & $\pi/2$ \\
\hline
$A'$ & & $+$ & 0 & $-$ & \\
\hline
$A$ & & $\nearrow$ & & $\searrow$ & 
\end{tabular}
\end{center}

$\left( \frac{\pi}{4} \cdot \tan \frac{\pi}{4} < 1 \text{ より } \frac{\pi}{4} < t_0 \text{ である。} \quad \cdots (*) \right)$

したがって
\[
A = 2 t_0 \cos t_0 = 2 \frac{\cos^2 t_0}{\sin t_0} \quad \cdots \text{①}
\]
となる。

次にBについて、$C'' = -\cos t \leqq 0 \ (-\frac{\pi}{2} \leqq t \leqq \frac{\pi}{2})$ より $y=\cos x$ のグラフは上に凸なので、Bは3頂点が $(\pm \pi/2, 0), (1, 0)$ の時の $\triangle$ の面積で、
\[
B = \frac{1}{2}\pi \quad \cdots \text{②}
\]

最後にCについて、中心は原点である。半径 $r$ とすると、対称性から $0 \leqq t \leqq \pi/2$ で
\[
t^2 + \cos^2 t \geqq r^2 \quad \cdots \text{③}
\]
となれば良い。この左辺を $f(t)$ とする。
\[
f'(t) = 2t - 2C \cdot S = 2t - \sin 2t \geqq 0 \quad (\because 0 \leqq t \leqq \pi/2 \text{ で } t \geqq \sin t)
\]
より、$f$ は単調増加だから、③を満たす $\max r$ は $r = \sqrt{0 + \cos 0} = 1$ となり、
\[
C = \frac{1}{2} \cdot 1^2 \cdot \pi = \frac{1}{2}\pi \quad \cdots \text{④}
\]
である。

以上これらの大小を比べる。(*)及び①からAが区間内で単調減少であることから
\[
A < 2 \frac{\cos^2 \frac{\pi}{4}}{\sin \frac{\pi}{4}} = \sqrt{2} < \frac{1}{2}\pi \quad (\because 8 < \pi^2 \text{ より } 2\sqrt{2} < \pi)
\]
だから②, ④とあわせて
\[
A < B = C \quad (\cdots \text{以下、原稿がここで途切れている})
\]

\end{document}
